\documentclass[10pt,a4paper]{amsart}
\usepackage[margin=19mm]{geometry}
\usepackage{amsmath,amssymb,mathtools}
\usepackage[scr=rsfs,cal=euler]{mathalfa}
\usepackage{xfrac}
\usepackage[unicode,hidelinks,bookmarks=false]{hyperref}
\newcommand{\ie}{i.e.\ }
\newcommand{\cf}{cf.\ }
\newcommand{\resp}{respectively\ }
\newcommand{\Subsection}{\subsection}
\providecommand{\nobreakdash}{\nobreak\hbox{-}}
\setlength{\emergencystretch}{2em}
\multlinegap0pt
\usepackage{amssymb}
\usepackage{enumerate}
\newtheorem{lemma}{Lemma}
\usepackage{cleveref}
\crefname{lemma}{Lemma}{Lemmas}
\newcommand{\Loc}{\operatorname{Loc}}
\newcommand{\cat}[1]{{\mathsf{#1}}}
\newcommand{\dcat}[2][]{\cat{D}_{#1}(#2)}
\newcommand{\perf}[1]{\mathsf{Perf}\,#1}
\newcommand{\thick}{\operatorname{thick}}
\title{High Frobenius pushforwards generate the\\ bounded derived category}
\author{Matthew R. Ballard, Srikanth B. Iyengar, Pat Lank, Alapan Mukhopadhyay, Josh Pollitz}
\date{Source: arXiv:2303.18085v3 (CC BY 4.0)}
\begin{document}
\maketitle

\noindent\emph{Source and licence.} High Frobenius pushforwards generate the\\ bounded derived category. Version arXiv:2303.18085v3 (CC BY 4.0), distributed by arXiv under the Creative Commons Attribution 4.0 International licence, http://creativecommons.org/licenses/by/4.0/, as recorded in the arXiv licence metadata for this exact version and shown on the abstract page https://arxiv.org/abs/2303.18085v3. Full licence notice: \texttt{licenses/arxiv-2303.18085-CC-BY-4.0.txt}.\par\smallskip

\noindent\textit{Editorial scope.} Counted unit: the article's lemma le:thick-Frobenius (source0.tex lines 848--855). The article's complete original proof of it is delivered with it (source0.tex lines 857--875, one \texttt{proof} environment of 19 lines). No mathematics is written, repaired or re-derived: the statement and the proof are the article's own lines, in the article's own notation, and no retained line is substituted or re-flowed.  Retained uncounted as the source-local context the proof needs: lines 342--345.  Nothing was omitted from any retained span, so the package carries no omission record.  No retained line contains a citation, so the wrapper carries no bibliography and no bibliography entry.  No retained line contains a figure environment, an image-inclusion command or drawing code, so no image is delivered and the package declares no asset dependency.  Editorial formatting only, in addition: an AMS \texttt{amsart} preamble that restates the article's own declarations and macros used by the retained lines, namely the environments \texttt{lemma} on separate counters, together with the standard packages those lines need.  The retained environments print in this document as Lemma 1 in order of appearance, whereas the article prints the same items with the numbering of its own full document.  The complete native source of the article as distributed in the arXiv e-print for this exact version is bundled unchanged as \texttt{sources/139516/source0.tex}.
\begin{equation}
\label{eq:neeman-cgt}
\Loc_{\cat K}(G) \cap {\cat K}^{\mathrm c}= \thick_{\cat K}(G)\,.
\end{equation}

\begin{lemma}
\label{le:thick-Frobenius}
Let $X$ be a noetherian scheme over a field of prime characteristic $p$. Let $U\subseteq \perf X$ be such that $\thick_{\dcat X}(U)=\perf X$. For any $G$ in $\dcat X$ and integer $e\ge 0$ one has
\[
\thick^{\odot}_{\dcat X}(F ^e_*G) 
 = \thick_{\dcat X}(F ^e_*(E\otimes G)\mid E\in U)\,.
\]
\end{lemma}

\begin{proof}
It suffices to verify this when $e=1$; this eases up the notation a bit. 

Let $F^*\colon \dcat{X}\to \dcat{X}$ denote pullback along $F$. For any collection of objects $C$ in $\dcat X$ with $\Loc_{\dcat X}(C)=\dcat{X}$, adjunction yields $\Loc_{\dcat X}(F^*C)=\dcat{X}$. From this observation and the fact that $F^*$ restricts to an endofunctor on $\perf{X}$, one can apply \cref{eq:neeman-cgt} to conclude that 
\begin{equation}\label{eq:pullback}
\thick_{\dcat X}(F^*U)=\thick_{\dcat X}(U)=\perf{X}\,.
\end{equation}
Now it remains to observe that 
\begin{align*}
 \thick_{\dcat X}(F _*(E\otimes G)\mid E\in U)&=\thick_{\dcat X}(F _*(F^*(E)\otimes G)\mid E\in U)\\
 &=\thick_{\dcat X}(E\otimes F _*(G)\mid E\in U)\\
 &=\thick^{\odot}_{\dcat X}(F_*G)\,;
\end{align*}
where the first equality is by \cref{eq:pullback}, the second equality is the projection formula
\[
F _*(F ^*(-)\otimes G)\simeq (-)\otimes F _*G
\]
on $\perf{X}$, and the last equality is evident since $\thick_{\dcat X}(U)=\perf{X}$. 
\end{proof}

\end{document}
